\documentclass[11pt]{article}
\usepackage[margin=1in]{geometry}
\usepackage{amsmath,amssymb,amsthm}
\usepackage{booktabs}
\usepackage{enumitem}
\usepackage{graphicx}
\usepackage{hyperref}

\newcommand{\dd}{\mathrm{d}}
\newcommand{\norm}[1]{\lVert #1 \rVert}
\newcommand{\code}[1]{\texttt{#1}}
\newcommand{\bs}[1]{\boldsymbol{#1}}
\newcommand{\omegah}{Omega\_h}
\newtheorem{proposition}{Proposition}
\newtheorem{remark}{Remark}
\emergencystretch=3em

\title{Coupled Energy-Minimization Smoothing\\
and Topological Mesh Adaptation:\\
Design of an Adaptivity Loop for Norma}
\author{Alejandro Mota}
\date{August 2026, revised September 2026}

\begin{document}
\maketitle

\begin{abstract}
Energy-minimization smoothing (EMS) improves a mesh by minimizing a
hyperelastic distortion energy over the nodal positions, with each element
measured against an ideal reference element. Smoothing alone, however,
terminates at configurations where no continuous motion of the nodes can
further reduce the energy, while the mesh remains poor because the
obstruction is topological. Numerical experiments show that alternating
EMS with targeted topological operations---edge swaps, collapses, and
splits applied to the elements of highest distortion energy---improves
mesh quality far beyond what either mechanism achieves in isolation. This
note designs such a coupled loop for Norma, informed by a detailed reading
of the adaptation logic of the \omegah{} library. Three questions are
addressed with concrete algorithms: when to trigger topological
operations, how to guarantee a monotonic decrease of the total distortion
energy across both smoothing and topological phases, and how to detect
stagnation of the smoothing phase with a scale-invariant criterion. The
central design decision is to use a single distortion functional, well
defined for any mesh topology, as the objective of both phases: with a
prescribed size or metric field, the metric-conditioned smoothing energy
now in Norma is such a functional; without one, an intrinsic
equal-volume functional takes its place. The revision of September 2026
fixes the scope (a preprocessing step on four-node tetrahedral meshes,
with no transfer of solution state), records the alternatives of wrapping
\omegah{} or keeping the hand-off to the Sierra tool \code{improve\_mesh},
and states why the topological operations should be accepted on the
energy with geometric quality as a floor. The note closes with a
step-by-step implementation plan, recommendations for topology and
adjacency data structures, and a ranked list of risks.
\end{abstract}

\tableofcontents

\section{Motivation and observed workflow}
\label{sec:motivation}

Mesh smoothing by energy minimization treats the mesh as a hyperelastic
body whose per-element reference configuration is an ideal element, and
minimizes the resulting stored energy over the nodal positions. The
energy acts as a barrier against element inversion and drives every
element toward the ideal shape. The mechanism is purely geometric: the
connectivity is fixed, and only the nodal positions move.

This fixity is also the limitation. Many poor meshes contain
configurations---slivers wedged between well-shaped neighbors, vertices
of excessive valence, edges shared by too many or too few
elements---whose distortion cannot be removed by any continuous motion of
the nodes. Smoothing then stalls: the energy plateaus, or the line search
fails against the inversion barrier while the gradient remains large. The
obstruction is topological, and the remedy is a local change of
connectivity. Conversely, topological operations applied to an unsmoothed
mesh are evaluated at poor nodal positions and are frequently rejected or
mis-ranked. The two mechanisms are complementary, and the literature on
tetrahedral mesh improvement confirms that their combination is
substantially stronger than either alone
\cite{freitag1997,klingner2007}.

The workflow this note develops is therefore iterative:
\begin{enumerate}[itemsep=2pt]
\item run EMS until the energy decrease stagnates;
\item identify the elements of highest distortion energy;
\item apply topological operations to the neighborhoods of those
  elements, accepting only operations that decrease the energy;
\item resume EMS, which the changed topology now permits to relax
  further;
\item repeat until neither mechanism makes progress.
\end{enumerate}
\paragraph{Scope.} The loop is a preprocessing step: it produces a mesh
before any analysis is run on it, so no solution state exists on the
mesh and no state transfer is involved. It targets four-node tetrahedra
only, which is the element type the smoothing supports and the type the
topological operations are written for; a mesh of quadratic tetrahedra
is obtained afterwards by conversion, placing the midside nodes at the
edge midpoints and projecting those on the boundary to the surface with
the level-set machinery already in place. The workflow it replaces is
staggered: Norma smooths, the smoothed mesh is passed to the Sierra tool
\code{improve\_mesh}, which changes the topology against a geometric
quality measure, and the result may be smoothed again
(Section~\ref{sec:alternatives}).

The design questions are precisely where the coupling is delicate: the
trigger for step 3, the acceptance criterion that keeps the total energy
monotone across steps 1--4, and the stagnation test in step 1. Sections
\ref{sec:triggers}--\ref{sec:stagnation} address them in turn, after three
preparatory sections: an inventory of the machinery Norma already
possesses (Section~\ref{sec:norma}), a summary of the decision logic
of \omegah{} as implemented in its source code
(Section~\ref{sec:omegah}), and an assessment of the alternatives to a
native implementation (Section~\ref{sec:alternatives}).

\section{Existing infrastructure in Norma}
\label{sec:norma}

\subsection{EMS as a solid-mechanics solve}

Norma implements EMS not as a separate module but as a mode of the
ordinary quasi-static \code{SolidMechanics} solve. A single substitution
in the element loop (\code{src/model.jl}) replaces each element's
reference configuration by an ideal regular tetrahedron before computing
the deformation gradient; everything else---assembly, solvers, boundary
conditions, output---is the standard machinery. The energy is the
Seth--Hill family with exponents $m=n=2$, which diverges both as the
element Jacobian tends to zero and as it tends to infinity, and therefore
acts as an inversion barrier. The nodal positions are the unknowns; the
minimization runs under the quasi-static integrator as pseudo-time
continuation, with matrix-free steepest-descent or L-BFGS steps.

The ideal reference element is built by \code{create\_smooth\_reference}
(restricted at present to four-node tetrahedra) as a regular tetrahedron
of edge length $h$, where $h$ follows one of four rules: equal volume to
the original element, average edge length of the original element, the
maximum of these two, or a user-supplied size field $s(t,x,y,z)$. A fifth
rule, \code{metric field} (September 2026), replaces the scalar size by a
metric tensor with three principal sizes and a rotation; the ideal
element is then anisotropic and the energy is evaluated on the
deformation gradient measured in the metric, so that the orientation of
the elements enters the objective (\code{docs/notes/ems-anisotropic}).
In all rules the input geometry is the \emph{original} mesh as read from
disk; the consequences of this choice for adaptivity are the subject of
Section~\ref{sec:intrinsic}.

Boundary nodes move but remain on the boundary, through either
axis-aligned Dirichlet conditions or the \code{Surface} boundary
condition, which represents a boundary surface as an analytic level set
$g(\bs{x})=0$ and enforces it by projection: nodes are returned to the
surface before assembly (\code{return\_to\_surface!}) and the gradient is
projected onto the tangent space afterward. This machinery is directly
reusable for placing and constraining nodes during topological
operations on curved boundaries.

\subsection{Assets already in place}

Three existing components carry much of the weight of the proposed loop.

\begin{enumerate}[itemsep=2pt]
\item \emph{A per-element distortion field.} The integrated element
  energy is stored per block and element in
  \code{model.stored\_energy}, updated on every evaluation of the model
  and already written to the Exodus output as an element variable. Step 2 of the workflow---finding
  the most distorted elements---requires no new computation.
\item \emph{Cross-mesh state transfer.} The module
  \code{src/field\_transfer.jl} performs a consistent $L^2$ transfer of
  nodal fields between meshes of the same geometry, and composes it
  with the variational recovery of \cite{mota2013} to remap
  integration-point internal variables. The preprocessing loop of this
  note does not need it; it is what a future in-simulation adaptivity
  capability would require, and it connects to the constrained
  projection of the companion note on internal-variable recovery in
  admissible spaces (\code{docs/notes/internal-variable-projection}).
\item \emph{Mid-run mesh replacement.} The module \code{src/swap.jl}
  already replaces the mesh of a running simulation: it rebuilds the
  model, integrator, and solver, and transfers the state to the new
  mesh. This is the orchestration pattern a topology change must follow,
  for a structural reason: \code{initialize\_storage} aliases the
  integrator and solver arrays into the model matrices through raw
  pointers (\code{unsafe\_wrap}), so any reallocation of the model
  arrays invalidates the integrator and solver. A topology change must
  therefore proceed by full rebuild, never by in-place resizing.
\end{enumerate}

\subsection{Gaps}

Equally important is what does not exist. The mesh is never held in
memory as a mutable object: the \code{ExodusDatabase} handle is the
source of truth, connectivity is re-read from the file on every access,
and no function writes coordinates or connectivity back. There are no
adjacency structures of any kind---no node-to-element, element-to-element,
or edge data---so there is at present no way to enumerate the cavity
around an edge. There is no element quality module; the de facto quality
measure of the code base is the smoothing energy itself. (Point
location, a linear scan when this note was first written, is now a
uniform-grid search in \code{src/spatial\_search.jl}.) These gaps
define Step 3 of the implementation plan in
Section~\ref{sec:architecture}.

\section{The adaptation logic of \omegah{}}
\label{sec:omegah}

\omegah{} \cite{omegah,ibanez2016} is a metric-driven simplex-adaptation
library whose decision logic is compact and, as verified directly against
its source code, organized around three ideas: a two-tier system of
thresholds, cavity operators filtered by acceptance criteria, and
progress-based termination. This section records the logic precisely,
because Section~\ref{sec:triggers} transplants it.

\subsection{Two tiers of thresholds}

The options structure \code{AdaptOpts} carries \emph{desired} and
\emph{allowed} values with distinct roles: desired values select
candidates for work, allowed values are hard limits that no accepted
operation may violate. The defaults are: metric edge lengths desired
within $[1/\sqrt{2},\,\sqrt{2}]$, length allowed up to $2\sqrt{2}$;
element quality desired at least $0.30$ (three dimensions), quality
allowed at least $0.20$. There is no lower allowed length: short edges
are only ever ``desired'' away. The quality measure is the
metric-conditioned mean ratio, an intrinsic shape measure equal to one
for a metric-equilateral simplex and non-positive for a degenerate or
inverted one.

\subsection{Control flow}

The top-level \code{adapt()} first evaluates an ``already satisfactory''
predicate---minimum quality at least the desired quality \emph{and} all
metric edge lengths within the desired band---and returns without
touching the mesh if it holds. Otherwise it runs two phases.

\begin{enumerate}[itemsep=2pt]
\item \emph{Length phase.} Alternate refinement of edges longer than
  $\sqrt{2}$ and collapse of edges shorter than $1/\sqrt{2}$, repeating
  until neither operator changes anything. Acceptance in this phase
  enforces only the hard limits: a refinement is accepted if the worst new
  element stays above the allowed quality; a collapse additionally may
  not create edges beyond the allowed length. Quality may therefore
  \emph{decrease} in this phase, down to the floor.
\item \emph{Quality phase.} While the minimum quality is below the
  desired value: attempt edge swaps; only if a full swap pass changes
  nothing, attempt sliver collapses; if both fail in the same pass,
  report failure and stop. Acceptance in this phase is strict
  improvement: a swap or sliver collapse is kept only if the minimum
  quality of the new cavity \emph{strictly exceeds} that of the old
  cavity. This strict filter is the only monotonicity guarantee in the
  library.
\end{enumerate}

Candidates for the quality phase are the elements below desired quality
\emph{dilated by four layers of adjacent elements}, because the operation
that repairs a bad element frequently lies on a neighboring entity.
Boundary edges are never swapped. Conflicts between overlapping cavities
are resolved by a Luby-style independent set in which the cavity of
higher quality wins, with a deterministic global-identifier tiebreak;
candidates not selected simply wait for the next pass.

\subsection{Termination}

There is no iteration counter anywhere in \code{adapt()}. Every operator
returns false exactly when no candidate survives its filter chain, and
each phase loop exits when its operators return false. Termination is
thus purely progress-based: the loop runs precisely as long as some
operation can be accepted. When the quality goal is unreachable, the
library says so and returns the mesh in its best achieved state rather
than looping.

\section{The Sierra hand-off and the alternatives to a native loop}
\label{sec:alternatives}

Three routes exist: keep the hand-off to \code{improve\_mesh}, wrap
\omegah{} into Norma, or implement the operators natively. The facts
about the two external tools, read from their sources in September 2026,
are these.

\paragraph{\code{improve\_mesh} and \code{emend}.} \code{improve\_mesh}
is a thin driver (about $2{,}500$ lines: an input parser, STK mesh
creation, and optional CAD geometry through \code{m2g}) over the Sierra
library \code{emend} (93 files, about $10{,}000$ lines). \code{emend}
implements edge swaps, face swaps, edge collapses (with a variant for
mandatory node removal), edge refinement, independent-set selection of
non-conflicting cavities, control of the deviation from the geometry
(\code{edgeGeomDeviationTol}, \code{sharpFeatureAngleTolInDegrees}), and
midnode straightening for quadratic elements. Its parameters follow the
\omegah{} logic: desired and allowed relative edge lengths at $1/\sqrt2$
and $\sqrt2$, a quality below which improvement is attempted,
iteration caps per operator, and termination when the quality stops
improving. It accepts nodal isotropic and anisotropic size fields,
Cartesian-aligned and cylindrical size fields, and a displacement field.
Its quality is an abstract \code{QualityMetric} compared through
\code{is\_first\_quality\_metric\_better\_than\_second}, so the measure
is pluggable, but acceptance is always the comparison of the worst
element of a cavity before and after the operation. It contains no field
transfer and no code from \omegah{}: it is an independent implementation
of the same decision logic on STK, and it is not open source.

\paragraph{\omegah{}.} C++14 with CMake; optional MPI, OpenMP, or CUDA;
tetrahedra and triangles; anisotropic metric fields with gradation
control; refinement, coarsening, swapping, and sliver coarsening under
\code{AdaptOpts}; Exodus input and output through SEACAS; boundary
geometry through EGADS or through entity classification alone. The
quality is the metric-conditioned mean ratio, fixed in the library; the
one extension hook, \code{UserTransfer}, concerns fields, not the
acceptance criterion. There is no C interface.

\paragraph{Comparison.} Table~\ref{tab:alternatives} sets the two
routes that would bring the operations into Norma against each other.

\begin{table}[ht]
\centering
\small
\begin{tabular}{p{2.7cm}p{5.4cm}p{5.4cm}}
\toprule
& Wrap \omegah{} & Native operators in Julia \\
\midrule
Work to first result &
  A C shim (mesh arrays in and out, metric per vertex, options, the
  \code{adapt} call) of a few hundred lines plus a \code{ccall} layer,
  packaged as an optional extension; about two weeks &
  Steps 3 to 7 of Section~\ref{sec:architecture}; three to four
  thousand lines of Julia with tests; two to three months to a validated
  loop \\
Build and distribution &
  A C++ dependency with CMake in a code that runs wherever Julia runs;
  CI on Ubuntu and macOS must build it; a binary recipe for users &
  None \\
Acceptance criterion &
  Fixed: metric-conditioned mean ratio and the length band; an energy
  criterion requires patching the library &
  Any: the energy criterion of Section~\ref{sec:monotonicity}, a geometric
  floor, or both \\
Coupling with smoothing &
  Loose: the two tools optimize different objectives &
  Tight: one objective, alternating phases, monotone descent \\
Boundary geometry &
  Classification only, or EGADS; the analytic \code{Surface} level sets
  would have to be re-expressed &
  \code{Surface} level sets and node and side sets reused directly \\
Parallelism &
  MPI and threads &
  Serial at first; threads later through independent sets \\
Maturity &
  High, deterministic, with stated quality guarantees &
  To be earned: cavity validity, orientation, and boundary preservation
  need careful tests \\
Public repository &
  Yes (BSD license) &
  Yes \\
\bottomrule
\end{tabular}
\caption{Bringing topological operations into Norma: wrapping \omegah{}
against a native implementation.}
\label{tab:alternatives}
\end{table}

\paragraph{Assessment.} Keeping the hand-off is mature and in use, but
it cannot be part of a public Norma workflow, and the two tools will
always optimize different quantities (the energy in Norma, a geometric
quality in \code{emend}), which is the source of the inefficiency of
the staggered loop. Wrapping \omegah{} buys a mature adapter quickly but
reproduces that staggered loop with a different external tool and adds
a build dependency. The value of doing this inside Norma is one
objective for both phases, and that requires the native route. The
native route is a bounded piece of work because its hard parts exist or
are specified: the energy kernel and the metric-space ideal element are
in the code, the geometry constraints exist, the orchestration pattern
exists, and the operators for four-node tetrahedra are well documented
(an edge loop of at most seven elements admits a small table of
triangulations; collapses and splits are local). \code{emend} is ten
thousand lines because it also handles quadratic elements, CAD
geometry, STK, and parallel independent sets, none of which are in the
first version here. The recommendation is the native route, with
\omegah{} in one narrow role: its command-line tools read and write
Exodus, so a mesh and a metric can be run through it and its length and
quality histograms compared with ours at no integration cost, and with
\code{improve\_mesh} as the reference for the validation cases of
Section~\ref{sec:architecture}.

\section{A single distortion functional}
\label{sec:intrinsic}

\subsection{The obstruction}

The design problem that governs the entire coupling is that the
smoothing objective under the \code{equal volume}, \code{average edge
length}, and \code{max} rules is tied to the original mesh: the ideal
reference of each element is built from the geometry of that element in
the mesh as read from disk. After a topological operation, a new element
has no original counterpart, so the objective is undefined on it; and
even where defined, an energy computed against history-dependent targets
cannot be compared across a topology change, because the comparison would
be between different functionals. A monotonicity guarantee requires a
single functional, defined for every admissible pair of topology and
nodal positions, that both phases of the loop demonstrably decrease.
There are two such functionals, depending on whether a target is
prescribed.

\subsection{With a prescribed size or metric field}
\label{sec:prescribed}

When the target is a size field $h(\bs X)$ or a metric field
$\bs M(\bs X)$ (\code{docs/notes/ems-anisotropic}), the ideal element
of any tetrahedron is the metric sampled at its centroid, for a new
element exactly as for an old one. The smoothing energy
\begin{equation}
D(\mathcal{T}, \bs{x}) \;=\; \sum_{e \in \mathcal{T}}
V^e_{\bs X}\, \psi\!\left(\bs F_{U,e}\right),
\qquad
\bs F_{U,e} = \bs F_M \bs F_e \bs F_M^{-1},
\label{eq:metric-functional}
\end{equation}
with $V^e_{\bs X} = h_1 h_2 h_3 V_{\mathrm{unit}}$ the volume of the
ideal element and $\bs F_{U,e}$ the deformation gradient measured in
the metric, is therefore defined for every admissible pair
$(\mathcal{T}, \bs x)$ without reference to the history of the mesh.
It measures size, shape, and orientation against the target in one
number: it vanishes only for a unit element of the metric, it diverges at
inversion, and it reduces to the isotropic energy for equal sizes. This
is the functional the loop should use whenever a target is prescribed,
and it changes two things in what follows. First, an operation that
moves the mesh toward the target size is a decrease of the same
functional the smoother minimizes: a too-large element has a large
volumetric energy and its children, each nearer the target volume, have
less; so size-driven splits and collapses need no separate phase and no
exemption from the acceptance test, and the length band of \omegah{}
survives only as the inexpensive way to select which edges to try.
Second, the volume floor of the rules \code{size field} and
\code{metric field}, which scales the target up to the volume of the
original element, exists only because smoothing cannot change the number
of elements; in the coupled loop it is unnecessary and the
\code{unrestricted} variants are the right ones. The target is held
fixed during a smoothing phase, sampled at the centroids of the elements
at the start of the phase, and re-sampled at every phase boundary, as
the anisotropic note describes; a new element created by a topological
operation receives the metric at its centroid when it is created. For
thresholds the natural density is $\hat w_e = \psi(\bs F_{U,e})$, the
energy per unit ideal volume, which is dimensionless with unit moduli.

\subsection{Without a prescribed target: the equal-volume functional}

Without a target, the mesh should keep its local size and only improve
its shape. Let $\mathcal{T}$ denote the mesh topology and $\bs{x}$ the
nodal positions. Define
\begin{equation}
D(\mathcal{T}, \bs{x}) \;=\; \sum_{e \in \mathcal{T}} W_e(\bs{x}),
\qquad
W_e(\bs{x}) \;=\; \int_{\Omega_e^{\text{ideal}}}
\psi\!\left(\bs{F}_e\right) \dd V,
\label{eq:intrinsic}
\end{equation}
where $\psi$ is the Seth--Hill energy density already in use,
$\bs{F}_e$ is the deformation gradient of element $e$ measured from its
ideal reference, and---this is the decision---the ideal reference of each
element is the regular tetrahedron of \emph{equal volume to that
element's own current volume}. This is the existing \code{"equal
volume"} rule of \code{create\_smooth\_reference}, applied to the
current rather than the original geometry.

With this choice the energy of an element depends only on its own
present shape: it is invariant under uniform scaling, vanishes exactly
for a regular tetrahedron, and diverges at degeneracy. The functional
\eqref{eq:intrinsic} is therefore \emph{intrinsic}---a function of
$(\mathcal{T}, \bs{x})$ alone, with no dependence on mesh history---and
is well defined for any tetrahedral mesh whatsoever. It belongs to the
same class as the mean-ratio quality of \omegah{}, which is likewise a
scale-invariant shape measure of the current element; the difference is
that \eqref{eq:intrinsic} is an energy with an inversion barrier rather
than a bounded quality, which is precisely what makes it minimizable by
the existing solvers. For threshold purposes define the normalized
density
\begin{equation}
\hat{w}_e \;=\; \frac{W_e}{V_e},
\label{eq:density}
\end{equation}
with $V_e$ the current element volume, so that $\hat{w}_e = 0$ for an
ideal element and $\hat{w}_e \to \infty$ at degeneracy. Note the
orientation: minimum quality in \omegah{} corresponds to
\emph{maximum} energy density here.

\begin{remark}
Size-targeting is excluded from \eqref{eq:intrinsic} by construction:
this functional keeps whatever size the mesh has. Operations toward a
target size belong to the functional \eqref{eq:metric-functional} of
Section~\ref{sec:prescribed}, where they are decreases of the objective
and need no separate phase. The two functionals are not mixed within
one loop: a run either has a prescribed target or it does not.
\end{remark}

\subsection{Re-baselining}

Re-baselining concerns the equal-volume functional; with a prescribed
target the ideal element is already fixed by the field. The smoothing
solver requires a \emph{fixed} objective during a solve:
if the ideal reference tracked the current volume continuously, the
target would move with the unknowns. The resolution is re-baselining at
phase boundaries: at the start of each smoothing phase, set the
smoothing reference configuration to the current coordinates and build
the equal-volume ideal references from it. The phase objective then
coincides with $D(\mathcal{T}, \cdot)$ at the start of the phase and
remains a fixed, minimizable functional throughout it. The relation
between the frozen-target phase objective and the intrinsic functional
is that of a majorize--minimize surrogate; the safeguard of
Proposition~\ref{prop:descent} converts the practical descent of $D$
into a guaranteed one.

\section{Question 1: adaptivity triggers}
\label{sec:triggers}

The translation of the \omegah{} logic into the energy setting is
summarized in Table~\ref{tab:translation} and specified below.

\begin{table}[ht]
\centering
\begin{tabular}{ll}
\toprule
\omegah{} & Proposed EMS loop \\
\midrule
mean-ratio quality (max.\ $1$, degenerate $\leq 0$) &
  energy density $\hat{w}_e$ (min.\ $0$, degenerate $\to\infty$) \\
\code{min\_quality\_desired} (candidate selection) &
  $\hat{w}_{\text{des}}$: element is worth working on \\
\code{min\_quality\_allowed} (hard floor) &
  $\hat{w}_{\text{all}} > \hat{w}_{\text{des}}$: hard ceiling \\
metric length band $[1/\sqrt{2}, \sqrt{2}]$ &
  same band selects split and collapse candidates \\
\code{nsliver\_layers} $= 4$ adjacency layers &
  same: dilate candidate set by $n_{\text{layers}}$ rings \\
swap first, sliver-collapse as fallback & same order \\
strict cavity-quality improvement filter &
  strict cavity-energy decrease (Section~\ref{sec:monotonicity}) \\
progress-based termination, no counters & same \\
\bottomrule
\end{tabular}
\caption{Correspondence between the \omegah{} adaptation logic and its
proposed energy-based analogue. The orientation inverts: minimum
quality corresponds to maximum energy density.}
\label{tab:translation}
\end{table}

\begin{enumerate}[itemsep=2pt]
\item \emph{Trigger for the topology phase.} Enter the topology phase
  when the smoothing phase has stagnated in the sense of
  Section~\ref{sec:stagnation} \emph{and} $\max_e \hat{w}_e >
  \hat{w}_{\text{des}}$. This is the ``already satisfactory'' predicate
  of \omegah{} with quality replaced by energy density; with a
  prescribed target the energy density already contains the length
  clause, since an element of the wrong size has a large volumetric
  energy.
\item \emph{Candidate set.} All elements with $\hat{w}_e >
  \hat{w}_{\text{des}}$, dilated by $n_{\text{layers}}$ rings of element
  adjacency (default four, following \omegah{}); the operation
  candidates are the interior edges of this set. Boundary edges are
  excluded from swapping; they participate in collapses and splits only
  under the surface constraints of Section~\ref{sec:architecture}.
\item \emph{Order within a pass.} Swaps first; collapses only when a
  full swap pass accepts nothing. Swaps move no nodes and therefore
  compose trivially with smoothing; collapses remove resolution and are
  the fallback. Splits participate where a bisection lowers the cavity
  energy (Remark~\ref{rem:splits}); with a prescribed target, edges
  longer than $\sqrt2$ in the metric are the split candidates and edges
  shorter than $1/\sqrt2$ the collapse candidates, in addition to the
  edges of the high-energy cavities.
\item \emph{Outer termination.} The coupled loop ends when the smoothing
  phase reports stagnation or convergence and the subsequent topology
  phase accepts zero operations, or when smoothing converges with
  $\max_e \hat{w}_e \leq \hat{w}_{\text{des}}$. Termination is
  progress-based, as in \omegah{}; a large safety cap on outer
  iterations is retained for defense but is not the semantic
  terminator. When the goal is unreachable, the loop reports it and
  returns the best achieved mesh.
\end{enumerate}

\section{Question 2: energy monotonicity}
\label{sec:monotonicity}

The reading of the \omegah{} source localizes its monotonicity exactly:
strict improvement is guaranteed only where the strict filter is
applied---swaps and sliver collapses---while refinement and size-driven
collapse may degrade quality to the floor. Adapting the criterion to the
EMS framework therefore means applying a strict-decrease filter, stated
in energy, to \emph{every} operation of the quality phase, and pairing
it with a monotone smoothing phase. Both are stated against the
intrinsic functional \eqref{eq:intrinsic}.

\subsection{The acceptance test}

A topological operation replaces the elements $C_{\text{old}}$ of a
cavity by new elements $C_{\text{new}}$ at fixed nodal positions
$\bs{x}$ (for a split, the new node is placed and locally relaxed first;
see Remark~\ref{rem:relax}). The operation is accepted if and only if
\begin{equation}
\sum_{e \in C_{\text{new}}} W_e(\bs{x})
\;\leq\;
\sum_{e \in C_{\text{old}}} W_e(\bs{x}) \;-\; \delta,
\qquad\text{and}\qquad
\max_{e \in C_{\text{new}}} \hat{w}_e \;\leq\; \hat{w}_{\text{all}},
\label{eq:acceptance}
\end{equation}
with $\delta > 0$ a small relative margin (for example $10^{-8}$ times
the old cavity energy) that precludes cycles of operations accepted at
zero gain. The cavity \emph{sum} replaces the cavity minimum of
\omegah{} deliberately: the sum is the restriction of the global
functional to the cavity, so local acceptance implies global descent,
which is the property the loop needs. The ceiling
$\hat{w}_{\text{all}}$ plays the role of the allowed-quality floor;
element inversion is excluded automatically because the Seth--Hill
barrier sends $\hat{w}_e$ to infinity there.

\paragraph{Energy and geometric quality.} The energy is the right
quantity to accept an operation on, for three reasons. It is one
functional for both phases, so the alternation cannot cycle and
terminates (Proposition~\ref{prop:descent}); a geometric criterion for
the topology and an energy for the smoothing are two objectives, each
phase can undo the other's progress, and termination then rests on
iteration caps, which is what \code{emend} does with its
\code{maxOuterIterations}. It measures size, shape, and orientation at
once: the scaled Jacobian and the mean ratio measure shape only, need a
separate length band for size, and are blind to the orientation of an
element relative to an anisotropic target unless the element is first
mapped into the metric space. And it carries the inversion barrier. The
geometric measures keep two roles. The first is the hard floor: the
test compares cavity sums, and a proposal could lower the sum while
creating one element worse than the analysis code tolerates, so the
second condition of \eqref{eq:acceptance} may be stated, in place of or in
addition to $\hat w_{\text{all}}$, as the minimum scaled Jacobian or
mean ratio the solver requires; both are cheap, and the geometric one
is what the analysts already ask for, so reporting it keeps the results
comparable with those of \code{improve\_mesh}. The second is candidate
selection through the length band, as stated above. One property of the
energy to keep in view: it is not scale-free across elements of
different target sizes, because the element energy is the density times
the ideal volume. Thresholds and the floor are therefore on the density
$\hat w_e$, while the cavity comparison is on the sums, which weight
elements by their target volume as a global objective should.

\subsection{The smoothing phase}

The L-BFGS path of the smoothing solver already performs an Armijo
backtracking line search on the energy, so each smoothing phase is
monotone in its own frozen-target objective; the steepest-descent path
currently backtracks on the residual norm instead and should use the
energy merit for this loop. With re-baselining at phase start
(Section~\ref{sec:intrinsic}), the phase objective agrees with
$D(\mathcal{T}, \cdot)$ at the initial iterate. Because the phase
objective holds its targets fixed while $D$ is intrinsic, the phase is
completed by a safeguard: evaluate $D$ once at phase end, and roll the
phase back if $D$ increased. The safeguard costs one evaluation and is
expected to be inert in practice; its purpose is to turn practical
descent into provable descent.

\begin{proposition}[Global descent and finite termination of the
topological work]
\label{prop:descent}
Let the coupled loop consist of smoothing phases with the end-of-phase
safeguard and topological operations accepted by the test
\eqref{eq:acceptance}. Then the sequence of values of $D$ sampled at phase
boundaries and after each accepted operation is nonincreasing and
bounded below by zero, hence convergent; each accepted operation
decreases $D$ by at least $\delta$ times the corresponding cavity
energy at acceptance, so if $\delta$ is enforced against a fixed
positive reference energy the number of accepted topological operations
is finite.
\end{proposition}

\begin{proof}
Only cavity elements change in a topological operation, so
\eqref{eq:acceptance} implies $D$ decreases by at least $\delta$ across it.
The safeguard makes each smoothing phase nonincreasing in $D$ by
construction. A nonincreasing sequence bounded below converges, and a
uniform positive decrement per accepted operation bounds their number.
\end{proof}

The guarantee lives in exactly one place: a single acceptance routine
implementing \eqref{eq:acceptance}, called by every operator. No operator
carries its own criterion.

\begin{remark}[Splits participate honestly]
\label{rem:splits}
Because the equal-volume energy is scale-invariant, bisecting a badly
stretched element into two better-shaped children genuinely lowers the
cavity energy, and refinement needs no exemption from the test
\eqref{eq:acceptance}. With a prescribed target the same holds for
size-driven refinement and coarsening, since the metric-conditioned
energy \eqref{eq:metric-functional} penalizes the wrong size directly.
\end{remark}

\begin{remark}[Local relaxation before the test]
\label{rem:relax}
For a split, place the new node at the edge midpoint---projected to the
surface by \code{return\_to\_surface!} when the edge lies on the
boundary---and run a few iterations of cavity-local smoothing with the
cavity boundary fixed before evaluating \eqref{eq:acceptance}. The cost is a
handful of element evaluations; the benefit is that each operation is
judged at its locally relaxed optimum rather than at an arbitrary
initial placement, which raises the acceptance rate substantially. This
is the natural advantage of an energy framework over the position-free
acceptance tests of \omegah{}, and it is the precise mechanism by which
the alternation outperforms either method alone: each accepted
operation unlocks relaxation that the previous topology forbade.
\end{remark}

\section{Question 3: stagnation detection}
\label{sec:stagnation}

The smoothing solver currently stops on a gradient-norm test with
absolute and relative tolerances---a \emph{convergence} criterion.
Stagnation, the absence of significant further decrease, requires a
criterion that is scale-free, tolerant of plateaus, and inexpensive.
The smoothing phase should end at the first of the following three
events, and report which one occurred.

\begin{enumerate}[itemsep=2pt]
\item \emph{Windowed relative decrease.} Let $W_n$ denote the phase
  objective at iteration $n$ and let $m$ be a window length (default
  $m = 10$). Stop when
  \begin{equation}
  W_{n-m} - W_n \;\leq\; \tau_{\text{stag}}\,
  \bigl(W_0 - W_n\bigr),
  \qquad \tau_{\text{stag}} = 10^{-3} \text{ by default.}
  \label{eq:stagnation}
  \end{equation}
  In words: the decrease achieved over the last $m$ iterations is a
  negligible fraction of the decrease achieved since the phase began.
  The test is invariant under affine rescaling of the energy and under
  problem size, and the window absorbs the short plateaus that
  quasi-Newton iterations produce near barrier-dominated
  configurations. Since Armijo backtracking makes $\{W_n\}$
  monotone, the left side is nonnegative and the test is well posed.
\item \emph{Convergence proper.} The existing projected-gradient test
  passes: the phase ended because the mesh is smooth, not because it is
  obstructed.
\item \emph{Topological lock.} The line search backtracks to failure,
  or the non-positive-Jacobian guard trips repeatedly, while the
  projected gradient remains large. This is not convergence: it is the
  inversion barrier certifying that no continuous nodal motion can
  improve the mesh from the current configuration. It is the strongest
  possible trigger for the topology phase and should be detected and
  reported explicitly.
\end{enumerate}

The soundness of \eqref{eq:stagnation} rests on the standard descent
inequality: for a continuously differentiable objective with Armijo
backtracking, the per-iteration decrease is bounded below by a multiple
of the squared projected-gradient norm times the accepted step
\cite{nocedal2006}. Sustained failure to decrease therefore certifies
either a vanishing projected gradient (event 2 in disguise) or a
collapsing step length against the barrier (event 3); the windowed test
detects both with a single properly scaled quantity and no
problem-dependent tuning.

Distinguishing the three events gives the outer loop its terminator at
no additional cost: event 2 with $\max_e \hat{w}_e \leq
\hat{w}_{\text{des}}$ means the loop is finished; events 1 and 3 call
for a topology phase; a topology phase that accepts nothing after
event 1 or 3 means the loop has reached the best this machinery can
achieve, which is reported rather than concealed---the exact analogue of
the explicit failure report of \omegah{} when its quality goal is
unreachable.

\section{Architectural plan}
\label{sec:architecture}

The implementation follows the repository convention of paired files:
\code{src/adapt\_types.jl} for the option and cache types,
\code{src/adapt.jl} for the methods. The steps are ordered so that each
is independently testable, with the first coupled loop appearing at
Step 4.

\begin{enumerate}[itemsep=3pt]
\item \emph{Energy kernel on subsets.} Factor the per-element energy
  evaluation out of the assembly loop into a kernel callable on an
  arbitrary subset of elements at given nodal positions, with the
  metric-sampled ideal element of Section~\ref{sec:prescribed} for any
  element and the equal-volume re-baselined variant as the fallback,
  and register the density $\hat{w}_e$ as an Exodus element variable. Testable immediately
  against \code{model.stored\_energy} on the existing \code{examples/ems}
  cases.
\item \emph{Stagnation criterion.} Add the windowed test
  \eqref{eq:stagnation} and the lock detection to the solver stopping
  logic, returning a phase outcome (converged, stagnated, locked)
  rather than a Boolean, and switch the steepest-descent merit to the
  energy. Independently testable: the strongly distorted examples
  should report the lock outcome where they currently exhaust
  iterations.
\item \emph{In-memory topology and adjacency.} A topology cache built
  once per topology phase from the Exodus-read connectivity:
  coordinates, per-block connectivity, and the three adjacencies the
  code base lacks---element-to-element across faces (by sorted-face
  matching), node-to-elements in compressed sparse row form, and an
  edge table mapping each edge to its incident elements. Boundary faces
  are identified as unpaired faces and cross-checked against the side
  sets; node-set and side-set membership is stored as per-entity flags
  so operations can update it. Four-node tetrahedra only, matching the
  existing restriction of \code{create\_smooth\_reference}.
\item \emph{Edge swap and the acceptance test; first coupled loop.}
  Implement the three-dimensional edge swap: extract the loop of
  elements around an edge (at most seven, following \omegah{}),
  enumerate the tabulated triangulations of the surrounding polygon,
  and select the best; implement the single acceptance routine
  \eqref{eq:acceptance}; wire the outer driver (smoothing phase, candidate
  collection, swap passes, repeat). Swaps move no nodes and need no
  state transfer, so this milestone alone should reproduce most of the
  observed benefit of the alternation.
\item \emph{Edge collapse.} Sliver removal under the acceptance test
  \eqref{eq:acceptance}; collapses on the boundary are restricted to
  boundary edges with the surviving node kept on the surface through
  the \code{Surface} level-set machinery; nodes belonging to node sets
  are never removed; an orientation guard rejects collapses that would
  invert any element (already implied by the ceiling, but checked
  cheaply before the energy evaluation).
\item \emph{Edge split and the size target.} Midpoint placement with
  surface projection for boundary edges, cavity-local relaxation per
  Remark~\ref{rem:relax}, then the acceptance test. With a prescribed target,
  edges longer than $\sqrt2$ or shorter than $1/\sqrt2$ in the metric
  join the candidate set, the target is used \code{unrestricted}, and a
  new element takes the metric at its centroid; these operations pass
  through the same test as the others. The acceptance case is the tube
  of the anisotropic note with a radial size of a quarter of the wall,
  which smoothing alone cannot reach.
\item \emph{Orchestration and output.} After each topology phase,
  rebuild the model, integrator, and solver through the existing
  mesh-replacement pattern of \code{src/swap.jl}, which avoids the
  pointer-aliasing hazard by construction. Write each adapted mesh as a
  fresh Exodus file in a numbered sequence: the Exodus format cannot
  change element counts within a file, so per-phase files are the
  practical route, and the replacement machinery already validates it.
  Input schema: an \code{adaptivity} block carrying
  $\hat{w}_{\text{des}}$, $\hat{w}_{\text{all}}$, the stagnation window
  and tolerance, the number of adjacency layers, and per-operator
  switches, mirroring \code{AdaptOpts}.
\item \emph{Validation against the hand-off.} For the cube, the
  tube, and the strongly distorted examples, run the same input mesh
  and target through \code{improve\_mesh} and through the native loop
  and compare the energy and scaled-Jacobian histograms of the two
  outputs; run \omegah{} on the same Exodus files where a metric is
  prescribed and compare its length and quality histograms. The first
  milestone at which this comparison is meaningful is Step 4, since
  swaps alone, alternated with the existing smoother, are expected to
  show most of the benefit.
\item \emph{Out of scope.} In-simulation adaptivity, which would invoke
  the driver from the evolution loop and transfer the solution state
  through \code{field\_transfer.jl}, and quadratic elements, which are
  obtained by conversion after the loop. Both are noted here so that the
  data structures of Section~\ref{sec:datastructures} do not preclude
  them.
\end{enumerate}

\section{Data structures and design patterns}
\label{sec:datastructures}

\begin{enumerate}[itemsep=3pt]
\item \emph{Rebuild per phase, not incremental maintenance.} \omegah{}
  rebuilds its topology arrays wholesale after each operator pass
  rather than patching them. The same pattern is right here: topology
  phases are rare relative to smoothing iterations, wholesale rebuild
  is simple and cache-friendly, and it eliminates the entire class of
  stale-adjacency defects. Within a pass, operations are applied
  through tombstones---dead elements and nodes are marked, new ones
  appended---and the arrays are compacted once at the end of the pass.
\item \emph{One cavity abstraction, one acceptance test.} Every operator produces
  a cavity proposal: the identifiers of the old elements, the proposed
  new connectivity, and optionally a new node. A single routine
  evaluates \eqref{eq:acceptance} and applies or discards the proposal. The
  monotonicity guarantee then has exactly one owner in the code.
\item \emph{Sequential priority order now, independent sets later.}
  The Luby-style independent set of \omegah{} exists for distributed
  parallelism. In serial, sorting the candidates by cavity energy in
  decreasing order and applying them sequentially is simpler and
  equally effective; a cavity that touches an element already modified
  in the current pass is deferred to the next pass, exactly as
  unselected \omegah{} candidates wait for theirs. If the topology
  phase is ever threaded, adopt the independent set with the same
  comparator: the better cavity wins, with a deterministic identifier
  tiebreak.
\item \emph{Exodus out of the hot path.} Connectivity is at present
  re-read from the database handle inside the assembly loop. The
  topology cache of Step 3 should become the in-memory source of truth
  during adaptation, with Exodus demoted to input and output at phase
  boundaries. This change benefits plain EMS performance even before
  any topological operation exists.
\item \emph{Outcome types over Booleans.} The solver reports converged,
  stagnated, or locked; each operator pass reports its number of
  accepted operations; the driver's control flow is written entirely
  against these outcomes. This makes the progress-based termination of
  Section~\ref{sec:triggers} auditable in the logs, in the same way
  that each \omegah{} operator returns false exactly when its filter
  chain leaves no candidate.
\end{enumerate}

\section{Risks and open questions}
\label{sec:risks}

Ranked by importance.

\begin{enumerate}[itemsep=2pt]
\item \emph{The acceptance test and the kernel.} The single functional of
  Section~\ref{sec:intrinsic} is the load-bearing design decision. It
  should be validated by a small standalone prototype---a single cavity
  with hand-verified energies before and after a swap, a collapse, and a
  split, with the metric-sampled ideal and with the equal-volume
  fallback---before Step 4 is undertaken.
\item \emph{Boundary collapses on curved surfaces.} The level-set
  projection and the energy test can conflict: a collapse that lowers
  the cavity energy may move the surviving node to a surface position
  that raises it again after projection. The test must therefore be
  evaluated \emph{after} projection, and boundary collapses should be
  expected to have a lower acceptance rate.
\item \emph{Four-node tetrahedra only.} This is a decision, not a
  restriction to be lifted: the smoothing and the operators are written
  for four-node tetrahedra, and higher-order meshes are produced by
  conversion after the loop. The input documentation should say so.
\item \emph{Energy thresholds across target sizes.} The element energy
  is the density times the ideal volume, so thresholds and the floor
  must be on the density and the cavity comparison on the sums
  (Section~\ref{sec:monotonicity}); a threshold on the element energy
  would favor large elements.
\item \emph{Geometric floor versus energy test.} An operation may lower
  the cavity energy while creating an element below the geometric floor
  the analysis requires; the floor is part of the acceptance test for that reason,
  and its value should be taken from the analysis code, not from the
  smoother.
\end{enumerate}

\section{Summary}
\label{sec:summary}

The coupled loop rests on one design decision and three algorithms. The
decision is to use a single distortion functional for both phases: the
metric-conditioned smoothing energy \eqref{eq:metric-functional} when a
size or metric field is prescribed, which is the common case for a
preprocessing step and makes size-driven operations decreases of the
same objective, and the intrinsic equal-volume functional
\eqref{eq:intrinsic} otherwise, so that in either case one quantity is
meaningful across every smoothing step and every topological operation.
The loop is a preprocessing step on four-node tetrahedra; it transfers
no solution state, and higher-order meshes are obtained by conversion
afterwards. Of the alternatives, wrapping \omegah{} would reproduce the
present staggered workflow with a new dependency, and the Sierra
hand-off cannot be part of a public workflow; both remain useful as
references for validation. The algorithms are: the transplanted \omegah{} trigger
logic---two-tier thresholds on the energy density, candidate dilation by
adjacency layers, swaps before collapses, progress-based
termination---; the acceptance test \eqref{eq:acceptance}, which strengthens
the strict-improvement filter of \omegah{} from the cavity minimum to
the cavity sum and thereby makes local acceptance imply global descent
(Proposition~\ref{prop:descent}); and the windowed stagnation test
\eqref{eq:stagnation} with explicit detection of the topological lock,
which supplies both the trigger for the topology phase and the
terminator of the outer loop. Norma already provides the per-element
energy field, the metric-space ideal element, the analytic surface
constraints, and the mesh-replacement orchestration; the genuinely new
construction is the in-memory topology cache with adjacency, the three
cavity operators, and the single acceptance test that all of them pass.

\begin{thebibliography}{9}

\bibitem{omegah}
D.~A.~Ibanez et al.
\omegah{}: reliable mesh adaptation.
Sandia National Laboratories,
\url{https://github.com/sandialabs/omega_h}.

\bibitem{ibanez2016}
D.~A.~Ibanez.
\emph{Conformal Mesh Adaptation on Heterogeneous Supercomputers}.
PhD thesis, Rensselaer Polytechnic Institute, 2016.

\bibitem{freitag1997}
L.~A.~Freitag, C.~Ollivier-Gooch.
Tetrahedral mesh improvement using swapping and smoothing.
\emph{International Journal for Numerical Methods in Engineering}
40 (1997) 3979--4002.

\bibitem{klingner2007}
B.~M.~Klingner, J.~R.~Shewchuk.
Aggressive tetrahedral mesh improvement.
In: \emph{Proceedings of the 16th International Meshing Roundtable},
Springer, 2007, 3--23.

\bibitem{nocedal2006}
J.~Nocedal, S.~J.~Wright.
\emph{Numerical Optimization}, 2nd edition.
Springer, 2006.

\bibitem{mota2013}
A.~Mota, W.~Sun, J.~T.~Ostien, J.~W.~Foulk~III, K.~N.~Long.
Lie-group interpolation and variational recovery for internal variables.
\emph{Computational Mechanics} 52 (2013) 1281--1299.

\end{thebibliography}

\end{document}
